% engines: lualatex, pdflatex
% expect-warnings: 0
% expect-errors: 0
\documentclass{article}
\usepackage[margin=1cm]{geometry}
\usepackage{../../latex/pgf-go}
\graphicspath{{../../doc/img/}}
\newcommand\case[1]{\par\medskip\noindent\textbf{#1}\par\nopagebreak}
\pagestyle{empty}
\begin{document}

% In PDF output a stone with an image is also drawn once and reused (a Form XObject), with its shadow
% and its shine, from its fifth use on; before that each look is drawn once in a TeX box and copied (a
% form used a few times makes the PDF heavier). Here the form comes at the second use. In every row the
% stone in the middle differs from the ones around it in one option only: if the look were reused by
% mistake, it would look like its neighbours.
\ExplSyntaxOn
\int_gset:Nn \g__pgfgo_stone_form_uses_int {1}
\ExplSyntaxOff

\case{Image: other file (C1), width 8pt (C2), rotate 45 (C3), opacity 0.4 (C4)}
\begin{goban}[5x4, profile = shell]
	\stones[b]{B1:4, D1:4}
	\stones[b, image = {file = pgfgo-yunzi-black.png}]{C1}
	\stones[b, image = {width = 8pt}]{C2}
	\stones[b, image = {file = pgfgo-baduktv-black.png, rotate = 45}]{C3}
	\stones[b, image = {opacity = 0.4}]{C4}
\end{goban}

\case{Shadow: no image (C1), silhouette (C2), offset 3pt (C3); shine image (C4), its silhouette (C5)}
\begin{goban}[5x5, profile = shell]
	\stones[b]{B1:5, D1:5}
	\stones[b, shadow = {image = {file = {}}}]{C1}
	\stones[b, shadow = {color = red, image = {silhouette}}]{C2}
	\stones[b, shadow = {offset = 3pt}]{C3}
	\stones[b, shine = {true, image = {file = pgfgo-shadow.png, width = 6pt}}]{C4}
	\stones[b, shine = {true, color = red, image = {file = pgfgo-shadow.png, width = 6pt, silhouette}}]{C5}
\end{goban}

% The silhouette is drawn after moving PGF's transformation to the PDF: its corners must still count in
% the size of the stone, or the form would cut the points of the stars.
\case{A silhouette bigger than the stone, turned: the whole stars are seen}
\begin{goban}[5x2]
	\stones[b, shadow = {true, color = blue, offset = 2pt,
		image = {file = pgfgo-star.pdf, width = 16pt, rotate = 18, silhouette}}]{B1, D2}
\end{goban}

% With several files or a random turn the image depends on the point: one \stones (left) and a \stone
% per point (right) must look the same, and so must the same points on a second goban.
\case{Image chosen and turned by position}
\begin{goban}[5x2, profile = shell]\stones[w]{A1, B1, C1, D1, E1, A2, C2, E2}\end{goban}\quad
\begin{goban}[5x2, profile = shell]
	\stone[w]{A1}\stone[w]{B1}\stone[w]{C1}\stone[w]{D1}\stone[w]{E1}\stone[w]{A2}\stone[w]{C2}\stone[w]{E2}
\end{goban}\quad
\begin{goban}[5x2]\stones[b, image = {file = pgfgo-baduktv-black.png, rotate = 0:90}]{A1, B1, C1, D1, E1}\end{goban}

\case{The same stones with scale = 0.6}
\begin{goban}[5x1, profile = shell, scale = 0.6]\stones[b]{B1, D1}\stones[w]{C1}\end{goban}

\newpage
\case{A second page reuses the looks of the first}
\begin{goban}[5x2, profile = shell]\stones[w]{A1, B1, C1, D1, E1}\stones[b]{A2, C2}\end{goban}

% With the usual threshold the first four stones of each look are copies of a box and the rest use the
% form: all the stones of a row, and the same point in every goban, must look the same. The translucent
% prisoners are drawn four times and then use a form with a transparency group.
\ExplSyntaxOn
\int_gset:Nn \g__pgfgo_stone_form_uses_int {5}
\ExplSyntaxOff
\case{Drawn, then a form: the same look seven times, the same points on six gobans}
\begin{goban}[7x1, profile = yunzi]\stones[b]{A1, B1, C1, D1, E1, F1, G1}\end{goban}

\begin{goban}[7x1, profile = yunzi]
	\prisoner[b]{A1}\prisoner[b]{B1}\prisoner[b]{C1}\prisoner[b]{D1}\prisoner[b]{E1}\prisoner[b]{F1}\prisoner[b]{G1}
\end{goban}

\begin{goban}[3x1, profile = yunzi]\stones[w]{A1, B1, C1}\end{goban}\quad
\begin{goban}[3x1, profile = yunzi]\stones[w]{A1, B1, C1}\end{goban}\quad
\begin{goban}[3x1, profile = yunzi]\stones[w]{A1, B1, C1}\end{goban}\quad
\begin{goban}[3x1, profile = yunzi]\stones[w]{A1, B1, C1}\end{goban}\quad
\begin{goban}[3x1, profile = yunzi]\stones[w]{A1, B1, C1}\end{goban}\quad
\begin{goban}[3x1, profile = yunzi]\stones[w]{A1, B1, C1}\end{goban}

\end{document}
